\section{Discussion}\label{sec:discussion}

STATADIC differs from conventional statistical interfaces mainly in the explicit
separation and traceability of scientific responsibilities. This is not a claim
that established libraries lack composability or validation: the agreement in
\cref{tab:agreement} rests on their careful conventions. STATADIC's emphasis is on
keeping lower-level components reusable while drawing a clear boundary between
software compatibility and scientific validity.

The approach has costs:
\begin{itemize}
  \item Explicit contracts require more design, documentation and review than
    convenience defaults, and the public interface is correspondingly more
    verbose. Roles such as exposure are declared on the model, and panel,
    spatial and survival workflows still require explicitly structured
    inputs.
  \item Formalization requires a faithful translation of published mathematics.
    A machine-checked proof of a mistranslated statement verifies nothing about
    the original result.
  \item Restricted persistence deliberately restores less than a pickled object
    would.
  \item Method-specific qualification is expensive and, as
    \cref{sec:evidence} shows, often adverse.
\end{itemize}

Several limitations apply at this revision:
\begin{itemize}
  \item Human review of the formal source mappings is pending, and no claim
    connects a Lean statement to compiled JAX execution.
  \item \texttt{Report.to\_latex} renders one model per tabular, so
    comparisons across software, such as \cref{tab:agreement}, are assembled
    by the analysis script.
  \item \texttt{PlotSpec} overlays series on shared axes but has no facets or
    logarithmic scales; \cref{fig:coverage} arranges four specifications in
    TeX.
  \item Regression prediction plots carry no uncertainty bands, because no
    qualified contract for regression prediction uncertainty exists yet.
  \item No performance claim extends beyond the measured CPU.
  \item The repository is access-restricted at the time of writing.
\end{itemize}

Formal foundations are also emerging in economics. MathEconLib
\citep{matheconlib-2026} is building a Lean library for mathematical economics,
with econometrics a long-term goal. Shared definitions could reduce the cost of
formalizing estimators. Integration should proceed through reviewed
correspondence between definitions, however, rather than by assuming that
independently developed libraries agree.

The longer-term value of the architecture is that mathematical claims,
executable methods, validation results and publication artifacts can be treated
as related but independently checkable research objects. An optional publishing
companion that assembles them into a verified evidence bundle has been designed
but not yet implemented.
